\section{Limits of sequences}

A sequence is a function whose input is usually a natural number. Therefore the most natural limiting question for a sequence is not what happens near a finite real input. The natural question is what happens when the index becomes larger and larger.

If \((a_n)\) is a sequence, we write
\[
\lim_{n\to\infty} a_n=L
\]
when the terms of the sequence approach the number \(L\) as \(n\) grows.

\begin{definitionbox}
We say that a sequence \((a_n)\) converges to \(L\) if the terms \(a_n\) become arbitrarily close to \(L\) for sufficiently large \(n\). We write
\[
\lim_{n\to\infty}a_n=L.
\]
If no finite limit exists, we say that the sequence diverges.
\end{definitionbox}

The phrase ``for sufficiently large \(n\)'' is important. A sequence may behave irregularly at the beginning and still have a limit. Limits describe long-term behavior, not the first few terms.

\subsection*{A first convergent sequence}

Consider
\[
a_n=\frac1n.
\]
The first terms are
\[
1,\quad \frac12,\quad \frac13,\quad \frac14,\quad \ldots
\]
These terms get closer and closer to \(0\). Therefore
\[
\lim_{n\to\infty}\frac1n=0.
\]

The sequence never actually reaches \(0\), because \(\frac1n>0\) for every natural number \(n\). But reaching the limit is not required. The question is whether the terms approach the limit.

\begin{remarkbox}
A limit is about approaching, not necessarily about becoming equal. The terms \(1/n\) approach \(0\), even though no term of the sequence is equal to \(0\).
\end{remarkbox}

\subsection*{Divergence to infinity}

Some sequences do not approach a finite number. For example,
\[
a_n=n
\]
grows without bound. We write informally
\[
\lim_{n\to\infty} n=\infty,
\]
meaning that the terms become larger than any fixed bound. This is not convergence to a real number. It is divergence to infinity.

\begin{examplebox}
The sequence
\[
a_n=n^2
\]
has terms
\[
1,4,9,16,\ldots
\]
The terms grow without bound. The sequence does not converge to a finite real number.
\end{examplebox}

\subsection*{Oscillation}

A sequence may also fail to converge because it oscillates.

\begin{examplebox}
Consider
\[
a_n=(-1)^n.
\]
The terms are
\[
-1,1,-1,1,\ldots
\]
They do not approach one number. The sequence alternates forever between two values. Therefore the sequence diverges.
\end{examplebox}

Oscillation does not automatically prevent convergence. If the oscillation becomes smaller, a limit may still exist.

\begin{examplebox}
Consider
\[
a_n=\frac{(-1)^n}{n}.
\]
The sign alternates, but the size of the terms becomes smaller:
\[
\left|\frac{(-1)^n}{n}\right|=\frac1n\to0.
\]
Therefore
\[
\lim_{n\to\infty}\frac{(-1)^n}{n}=0.
\]
\end{examplebox}

The difference between these two examples is the size of the oscillation. In \((-1)^n\), the oscillation never becomes small. In \((-1)^n/n\), the oscillation is squeezed toward \(0\).

\subsection*{Reading sequence limits}

When reading a sequence limit, we should identify the long-term structure of the formula. Does the numerator and denominator grow at comparable rates? Does a power dominate? Does a term go to zero? Does an alternating sign matter?

\begin{examplebox}
Consider
\[
a_n=\frac{n}{n+1}.
\]
For large \(n\), the numerator and denominator have almost the same size. Divide numerator and denominator by \(n\):
\[
\frac{n}{n+1}=\frac{1}{1+\frac1n}.
\]
Since \(\frac1n\to0\), we get
\[
\lim_{n\to\infty}\frac{n}{n+1}=1.
\]
\end{examplebox}

This calculation should be understood conceptually: the extra \(1\) in the denominator becomes insignificant compared with \(n\) when \(n\) is large.

\begin{summarybox}
When studying a sequence limit, ask:
\begin{itemize}
\item What happens when \(n\) becomes large?
\item Are the terms approaching a finite number?
\item Are the terms growing without bound?
\item Is there oscillation?
\item If there is oscillation, does its size shrink?
\item Which part of the formula dominates for large \(n\)?
\end{itemize}
\end{summarybox}

A sequence limit describes long-term behavior. The first few terms may help us guess, but the limit is determined by what happens as \(n\) grows without bound.